%-------------------------------------------------------------------------
%	PACKAGES AND OTHER DOCUMENT CONFIGURATIONS
%-------------------------------------------------------------------------
% Fix PDF version compatibility for including Front/main.pdf when pdfTeX is used.
\ifdefined\pdfminorversion
  \pdfminorversion=7
\fi

% Include pdf pages in the document
% Necessary to include the front pages (cover and etc.)
\usepackage{pdfpages}
% Restores the hyperlinks embedded in Front/main.pdf when the cover is imported.
% Its annotations are generated by Front/extract-links.tex with LuaLaTeX.
\usepackage{pax}

% Fix top page geometry on long titles
\setlength{\headheight}{14pt}  %Try fix error

% Language hyphenation and typographical rules
\usepackage[portuguese,english]{babel}
%Custom hyphenization
\hyphenation{Py-thon}
\hyphenation{Ju-py-ter}
\hyphenation{Ma-the-ma-ti-ca}

% Glossaries e Abreviações - deve vir DEPOIS do babel e hyperref
\usepackage[acronym,toc,nonumberlist]{glossaries}
\setacronymstyle{long-short}

% Inline quotes
% added for \begin{displayquote}
\usepackage[autostyle]{csquotes}

% Bibliography setup
% Use the natbib reference package - read up on this to edit the reference
% style; if you want text (e.g. Smith et al., 2012) for the in-text references
% (instead of numbers), remove 'numbers'
\usepackage[square, numbers, comma, sort&compress]{natbib}
%\bibliographystyle{IEEEtranN}  % I actually quite like this one
% \bibliographystyle{apsrev4-1-etal} % With emphasized titles. ORIGINAL
% Prevent that the first citation is in the ToC
\usepackage{notoccite}

% Interesting float placements (like 'H') and custom float types
\usepackage{float}
% Text wrapped around pictures
% https://pt.sharelatex.com/learn/Wrapping_text_around_figures
\usepackage{wrapfig}
% Force float barriers, use as \FloatBarrier
\usepackage[section]{placeins}
% Place floats *above* footnotes
\usepackage[bottom, perpage, symbol]{footmisc}
% Set default float placement
\makeatletter
\renewcommand{\fps@figure}{tbph}
\renewcommand{\fps@table}{tbph}
\makeatother

% Pretty colours
\usepackage{xcolor}
% \usepackage{color} % Deprecated by xcolor

% SVGs with Inkscape and PDF+LaTeX
% https://tex.stackexchange.com/questions/473994/svg-and-inkscape
\usepackage[inkscapearea=page]{svg}
% Specifies the directory where vector are stored
\svgpath{{Svgs/}}

% Graphics stuff
\usepackage{graphicx}  % invoked by svg
% Specifies the directory where pictures are stored
\graphicspath{{Figures/}}

% For sub-figures and stuff
\usepackage{caption}
\usepackage{subcaption}

% Math stuff
\usepackage{amsmath} % Interesting environments
\usepackage{amssymb} % Interesting symbols
\usepackage{commath} % Interesting macros
\usepackage{braket} % Dirac bra-ket and set notations
\usepackage{mathpazo} % Math font (palatino for Computer Modern on math)
\usepackage{mathtools} % Mathematical tools to use with amsmath
% Math alphabet
\DeclareMathAlphabet{\pazocal}{OMS}{zplm}{m}{n}
\newcommand{\Sa}{\pazocal{S}}
\newcommand{\Ua}{\pazocal{U}}
\newcommand{\Ha}{\pazocal{H}}
\newcommand{\Fa}{\pazocal{F}}
\newcommand{\Ia}{\pazocal{I}}
\newcommand{\Ea}{\pazocal{E}}
\newcommand{\ja}{\pazocal{J}}
%Custom math operators
\DeclareMathOperator*{\meshgrid}{meshgrid}
% Floor and ceiling of numbers
\DeclarePairedDelimiter\ceil{\lceil}{\rceil}
\DeclarePairedDelimiter\floor{\lfloor}{\rfloor}
% Notation variables
\newcommand{\dd}{\mathrm{d}}

% Units and numbers in text
\usepackage{siunitx}
\DeclareSIUnit\baud{Bd} % Baud

% Reimplementation of and extensions to LaTeX verbatim
\usepackage{verbatim} %added for \begin{comment}

% Fancy chapter start quotes
\usepackage{epigraph, varwidth}
\input{epigraphoverload}

% Use more than one optional parameter in a new commands
\usepackage{xargs}

% Code listings
\usepackage{listings}
% Colors for the listing
\definecolor{dkgreen}{rgb}{0,0.6,0}
\definecolor{gray}{rgb}{0.5,0.5,0.5}
\definecolor{mauve}{rgb}{0.58,0,0.82}
\definecolor{codegreen}{rgb}{0,0.6,0}
\definecolor{codegray}{rgb}{0.5,0.5,0.5}
\definecolor{codepurple}{rgb}{0.58,0,0.82}
\definecolor{backcolour}{rgb}{0.95,0.95,0.92}
\definecolor{orange}{RGB}{255,127,0}
% Python style for code blocks
\lstdefinestyle{Python}{
        language=Python,
         numberstyle=\small,
         stepnumber=2,
         numbersep=10pt,
         basicstyle={\small\ttfamily},
         keywordstyle    = \color{blue},
         commentstyle    = \color{red}\ttfamily,
         stringstyle=\color{orange},
         tabsize=2,
         columns=fullflexible,
         backgroundcolor=\color{backcolour},
         frame=none,
         numbers=left,
         aboveskip=5mm,
         belowskip=5mm,
         breaklines=true
}

% Algorithmicx provides a flexible, yet easy to use, way for inserting good
% looking pseudocode or source code in your papers.
\usepackage{algorithmicx}

% Hyperref and Backref
% backref makes the bibliography say where the entry was cited.
% For the print version of the thesis you might wanna set all colors to back
\usepackage{hyperref}
\usepackage[hyperpageref]{backref}
% Define a darker, less prominent blue for links
\definecolor{darkblue}{RGB}{0,51,102}
\hypersetup{colorlinks, citecolor=darkblue, urlcolor=darkblue,
        linkcolor=darkblue, breaklinks=true, hypertexnames=true}
\renewcommand*{\backref}[1]{}
\renewcommand*{\backrefalt}[4]{%
    \ifcase #1%
          \or [Cited on page~#2.]%
          \else [Cited on pages~#2.]%
    \fi%
    }
% Interesting URL breakings
\usepackage{url}
\def\UrlBreaks{\do\/\do-\do\&\do.\do:}

% Variants of \fbox and other games with boxes
\usepackage{fancybox}

% LaTeX default text is fully-justified, but often left-justified text may be a
% more suitable format. This left-alignment can be easily accomplished by
% importing the ragged2e package.
\usepackage{ragged2e}

% Create tabular cells spanning multiple rows
\usepackage{multirow}

% Changes bullet points marker
\renewcommand{\labelitemi}{\(\bullet\)}

% Notes on the documents
% https://tex.stackexchange.com/questions/9796/how-to-add-todo-notes
% https://tex.stackexchange.com/questions/316220/todo-commentsnot-include-and-left-align
% Examples:
% \unsure{Is this correct?}, \change{Change this!},
% \info{This can help me in chapter seven!}
% \improvement{This really needs to be improved!\\ What was I thinking?!}
% \thiswillnotshow{This is hidden since option `disable' is chosen!}
% WARNING: It eliminates whitespaces in front of it.
% You can add trailing {} to avoid.

\usepackage[colorinlistoftodos,
    prependcaption,
    textsize=tiny,
    textwidth=2cm]
        {todonotes}
% You can add:
% \setlength{\marginparwidth}{3cm}\reversemarginpar
% before \todo on each command for a different effect
\newcommandx{\unsure}[2][1=]{
    % \setlength{\marginparwidth}{3cm}\reversemarginpar
    \todo[linecolor=red,backgroundcolor=red!25,bordercolor=red,#1]{#2}
    }
\newcommandx{\change}[2][1=]{
    % \setlength{\marginparwidth}{3cm}\reversemarginpar
    \todo[linecolor=blue,backgroundcolor=blue!25,bordercolor=blue,#1]{#2}
    }
\newcommandx{\info}[2][1=]{
    % \setlength{\marginparwidth}{3cm}\reversemarginpar
    \todo[linecolor=green,backgroundcolor=green!25,bordercolor=green,#1]{#2}
    }
\newcommandx{\improvement}[2][1=]{
    % \setlength{\marginparwidth}{3cm}\reversemarginpar
    \todo[linecolor=yellow,backgroundcolor=yellow!25,bordercolor=yellow,#1]{#2}
    }
\newcommandx{\thiswillnotshow}[2][1=]{\todo[disable,#1]{#2}}

% Tradução dos nomes das listas para português
\addto\captionsportuguese{%
  \renewcommand{\listfigurename}{Lista de Figuras}%
  \renewcommand{\listtablename}{Lista de Tabelas}%
  \renewcommand{\contentsname}{Índice}%
}

% Comando para adicionar gráfico na Lista de Gráficos
% Use \graphiccaption{texto da legenda} em vez de \caption{texto da legenda}
% para que o gráfico apareça na Lista de Gráficos
\newcommand{\graphiccaption}[2][]{%
  \captionsetup{list=no}%
  \ifx\\#1\\
    \caption{#2}%
    \addcontentsline{logr}{figure}{\protect\numberline{\thefigure}#2}%
  \else
    \caption[#1]{#2}%
    \addcontentsline{logr}{figure}{\protect\numberline{\thefigure}#1}%
  \fi
}

% Inicializar glossário - DEVE vir no final do preamble
% Variant for a shared caption placed outside a figure environment.
\newcommand{\graphiccaptionof}[2][]{%
  \captionsetup{type=figure,list=no}%
  \graphiccaption[#1]{#2}%
}

\makeglossaries

% Remover espaço extra antes de acrônimos
\glssetnoexpandfield{first}

% Comando customizado para acrônimo curto com link (sem espaço extra)
\newcommand{\acr}[1]{\glsdisp{#1}{\glsentryshort{#1}}}

% Definir abreviações aqui (ou num arquivo separado)
\newacronym{fcup}{FCUP}{Faculdade de Ciências da Universidade do Porto}
\newacronym{up}{UP}{Universidade do Porto}
\newacronym{cnn}{CNN}{Convolutional Neural Network}
\newacronym{gan}{GAN}{Generative Adversarial Network}
\newacronym{dl}{DL}{Deep Learning}
\newacronym{unet}{U-Net}{U-shaped Network}
\newacronym{ct}{CT}{Computed Tomography}
\newacronym{hu}{HU}{Unidade Hounsfield}
\newacronym{mar}{MAR}{Metal Artifact Reduction}
\newacronym{nmar}{NMAR}{Normalized Metal Artifact Reduction}
\newacronym{fbp}{FBP}{Filtered Back Projection}
\newacronym{ir}{IR}{Iterative Reconstruction}
\newacronym{psnr}{PSNR}{Peak Signal-to-Noise Ratio} 
\newacronym{ssim}{SSIM}{Structural Similarity Index Measure}
\newacronym{ms-ssim}{MS-SSIM}{Multi-Scale Structural Similarity Index Measure}
\newacronym{mae}{MAE}{Mean Absolute Error}
\newacronym{dice}{DICE}{Dice Similarity Coefficient}
\newacronym{rmse}{RMSE}{Root Mean Square Error}
\newacronym{roi}{ROI}{Region of Interest}
\newacronym{bce}{BCE}{Binary Cross-Entropy}
\newacronym{tcia}{TCIA}{The Cancer Imaging Archive}
\newacronym{gpu}{GPU}{Graphics Processing Unit}
\newacronym{mbir}{MBIR}{Model-Based Iterative Reconstruction}
\newacronym{omar}{O-MAR}{Orthopedic Metal Artifact Reduction}
\newacronym{imar}{iMAR}{Iterative Metal Artifact Reduction}
\newacronym{semar}{SEMAR}{Single Energy Metal Artifact Reduction}
\newacronym{amp}{AMP}{Automatic Mixed Precision}
\newacronym{relu}{ReLU}{Rectified Linear Unit}
\newacronym{bn}{BN}{Batch Normalization}
\newacronym{npy}{NPY}{NumPy Array}
\newacronym{ctmar}{CT-MAR}{CT Metal Artifact Reduction Challenge}
